% GENERATED FILE. Edit canonical lessons, then run npm run generate:books.
% canonical-source-hash: fnv1a64:114ecc08a08723e9
% canonical-lessons: FR-C342-evaluer, FR-C342-juger, FR-C342-jugement, FR-C342-subjectif, FR-C342-impartial

\chapter{Évaluer et juger}
\label{ch:fr-evaluer-et-juger}
\hlchaptermodality{342}

\begin{chapteropening}
\textbf{By the end of this chapter:} \emph{I can evaluate something, judge it, give a judgement and say whether a view is subjective or impartial.}

Complete the last lesson of chapter 342: I can evaluate something, judge it, give a judgement and say whether a view is subjective or impartial.
\end{chapteropening}

\section[évaluer]{évaluer — to evaluate, to assess}

\label{lesson:FR-C342-evaluer}



\begin{quote}
Before the new one: say the French for the precision, then the French for to clear up.
\end{quote}

\subsection*{You'll want to know: évaluer}
\textbf{évaluer} — \textquotedblleft{}to evaluate, to assess\textquotedblright{}.

\begin{quote}
Le professeur évalue notre travail chaque mois.
\end{quote}

\begin{quote}
\emph{The teacher assesses our work every month.}
\end{quote}

\emph{Évaluer le prix d'une maison} — to estimate it.

\begin{cousinweb}
\emph{évaluer} is usually explained as a French formation from \emph{é-} and Old French \emph{value} \textquotedblleft{}worth\textquotedblright{}, the old participle of \emph{valoir}, which is inherited from Latin \emph{valēre} \textquotedblleft{}to be strong, to be worth\textquotedblright{} (Lewis \& Short).
\end{cousinweb}

\begin{grammarlens}[title={What you've built}]
One more word for telling, explaining and describing at B1.
\end{grammarlens}

\subsection*{Your turn}
\begin{itemize}
\raggedright
  \item \emph{Say it:} \emph{évaluer}
  \item \emph{Say it:} \emph{évaluer}, once more
  \item \emph{Say it:} a sentence of your own with \emph{évaluer}
  \item \emph{Recall:} say the French for the precision, then the French for to clear up, then say \emph{évaluer} again
\end{itemize}

\begin{tcolorbox}[breakable,colback=teal!4,colframe=teal!35!black,title={Before you move on}]
What does évaluer mean? (\textquotedblleft{}To evaluate, to assess\textquotedblright{}.) Say it once more.
\end{tcolorbox}
\section[juger]{juger — to judge}

\label{lesson:FR-C342-juger}



\begin{quote}
Before the new one: say the French for to clear up, then the French for to evaluate.
\end{quote}

\subsection*{You'll want to know: juger}
\textbf{juger} — \textquotedblleft{}to judge\textquotedblright{}.

\begin{quote}
Il ne faut pas juger les gens sur leurs vêtements.
\end{quote}

\begin{quote}
\emph{You should not judge people by their clothes.}
\end{quote}

\emph{Nous jugeons}, with an \emph{e} to keep the \emph{g} soft.

\begin{cousinweb}
\emph{juger} is inherited from Latin \emph{iūdicāre} \textquotedblleft{}to pass judgement\textquotedblright{} (EtymDB; Lewis \& Short), from \emph{iūdex} \textquotedblleft{}judge\textquotedblright{}, itself built from \emph{iūs} \textquotedblleft{}law\textquotedblright{} and \emph{dīcere} \textquotedblleft{}to say\textquotedblright{}.
\end{cousinweb}

\begin{grammarlens}[title={What you've built}]
One more word for telling, explaining and describing at B1.
\end{grammarlens}

\subsection*{Your turn}
\begin{itemize}
\raggedright
  \item \emph{Say it:} \emph{juger}
  \item \emph{Say it:} \emph{juger}, once more
  \item \emph{Say it:} a sentence of your own with \emph{juger}
  \item \emph{Recall:} say the French for to clear up, then the French for to evaluate, then say \emph{juger} again
\end{itemize}

\begin{tcolorbox}[breakable,colback=teal!4,colframe=teal!35!black,title={Before you move on}]
What does juger mean? (\textquotedblleft{}To judge\textquotedblright{}.) Say it once more.
\end{tcolorbox}
\section[le jugement]{le jugement — the judgement}

\label{lesson:FR-C342-jugement}



\begin{quote}
Before the new one: say the French for to evaluate, then the French for to judge.
\end{quote}

\subsection*{You'll want to know: le jugement}
\textbf{le jugement} — \textquotedblleft{}the judgement\textquotedblright{}.

\begin{quote}
J'ai confiance en ton jugement.
\end{quote}

\begin{quote}
\emph{I trust your judgement.}
\end{quote}

In court, \emph{un jugement} is a ruling.

\begin{cousinweb}
\emph{jugement} is inherited through Old French \emph{jugement} from Late Latin \emph{iūdicāmentum} (EtymDB), built on \emph{iūdicāre} \textquotedblleft{}to judge\textquotedblright{} and \emph{iūdex} \textquotedblleft{}judge\textquotedblright{} (Lewis \& Short).
\end{cousinweb}

\begin{grammarlens}[title={What you've built}]
One more word for telling, explaining and describing at B1.
\end{grammarlens}

\subsection*{Your turn}
\begin{itemize}
\raggedright
  \item \emph{Say it:} \emph{le jugement}
  \item \emph{Say it:} \emph{le jugement}, once more
  \item \emph{Say it:} a sentence of your own with \emph{le jugement}
  \item \emph{Recall:} say the French for to evaluate, then the French for to judge, then say \emph{le jugement} again
\end{itemize}

\begin{tcolorbox}[breakable,colback=teal!4,colframe=teal!35!black,title={Before you move on}]
What does le jugement mean? (\textquotedblleft{}The judgement\textquotedblright{}.) Say it once more.
\end{tcolorbox}
\section[subjectif]{subjectif — subjective}

\label{lesson:FR-C342-subjectif}



\begin{quote}
Before the new one: say the French for to judge, then the French for the judgement.
\end{quote}

\subsection*{You'll want to know: subjectif}
\textbf{subjectif} — \textquotedblleft{}subjective\textquotedblright{}.

\begin{quote}
C'est un avis très subjectif.
\end{quote}

\begin{quote}
\emph{It is a very subjective view.}
\end{quote}

Feminine \emph{subjective}: \emph{une impression subjective}.

\begin{cousinweb}
\emph{subjectif} comes from Latin \emph{subiectīvus} (EtymDB), formed on \emph{subiectus}, the participle of \emph{subicere} \textquotedblleft{}to throw under, to place under\textquotedblright{} (Lewis \& Short), built on \emph{iacere} \textquotedblleft{}to throw\textquotedblright{}.
\end{cousinweb}

\begin{grammarlens}[title={What you've built}]
One more word for telling, explaining and describing at B1.
\end{grammarlens}

\subsection*{Your turn}
\begin{itemize}
\raggedright
  \item \emph{Say it:} \emph{subjectif}
  \item \emph{Say it:} \emph{subjectif}, once more
  \item \emph{Say it:} a sentence of your own with \emph{subjectif}
  \item \emph{Recall:} say the French for to judge, then the French for the judgement, then say \emph{subjectif} again
\end{itemize}

\begin{tcolorbox}[breakable,colback=teal!4,colframe=teal!35!black,title={Before you move on}]
What does subjectif mean? (\textquotedblleft{}Subjective\textquotedblright{}.) Say it once more.
\end{tcolorbox}
\section[impartial]{impartial — impartial, fair}

\label{lesson:FR-C342-impartial}



\begin{quote}
Before the new one: say the French for the judgement, then the French for subjective.
\end{quote}

\subsection*{You'll want to know: impartial}
\textbf{impartial} — \textquotedblleft{}impartial, fair\textquotedblright{}.

\begin{quote}
Un juge doit rester impartial.
\end{quote}

\begin{quote}
\emph{A judge must stay impartial.}
\end{quote}

Feminine \emph{impartiale}, plural \emph{impartiaux}.

\begin{cousinweb}
\emph{impartial} is built from \emph{in-} \textquotedblleft{}not\textquotedblright{} and \emph{partial} \textquotedblleft{}taking sides\textquotedblright{}, a learned formation on Latin \emph{pars, partis} \textquotedblleft{}part, side\textquotedblright{} (Lewis \& Short).
\end{cousinweb}

\begin{grammarlens}[title={What you've built}]
One more word for telling, explaining and describing at B1.
\end{grammarlens}

\subsection*{Your turn}
\begin{itemize}
\raggedright
  \item \emph{Say it:} \emph{impartial}
  \item \emph{Say it:} \emph{impartial}, once more
  \item \emph{Say it:} a sentence of your own with \emph{impartial}
  \item \emph{Recall:} say the French for the judgement, then the French for subjective, then say \emph{impartial} again
\end{itemize}

\begin{tcolorbox}[breakable,colback=teal!4,colframe=teal!35!black,title={Before you move on}]
What does impartial mean? (\textquotedblleft{}Impartial, fair\textquotedblright{}.) Say it once more.
\end{tcolorbox}
